\documentclass{article}
\usepackage{amsmath}
\usepackage{amssymb}
\usepackage[margin=1in]{geometry}
\usepackage{enumitem}
\usepackage{graphicx}
\newcommand{\sectionbreak}{\clearpage}
\newcommand{\qed}{$\hfill \square$}


\title{Proof Workout Mastery Workbook}
\author{Joe Ricks}
\date{\today}

\begin{document}

\maketitle

\textbf{\Huge I have neither given nor received unauthorized assistance.}

\sectionbreak
\textbf{Warm Up THM 1} (Given Example)
\medbreak
\textit{Prove, for all positive integers, $A$, $B$, if $A$ and $B$ are even, then $C = A + B$ is even:}

\bigbreak

\begin{tabular}{r@{}l@{\hspace{2cm}}l}
\multicolumn{2}{l}{There is $d$ such that $A = 2d$} & definition of even \\[6pt]
\multicolumn{2}{l}{There is $h$ such that $B = 2h$} & definition of even (\textit{notice $d$/$h$ cannot be the same, why?}) \\[6pt]
$C$ & ${}= A + B$ & premise \\[6pt]
& ${}= 2d + 2h$ & substitution \\[6pt]
& ${}= 2(d + h)$ & factoring \\[6pt]
\multicolumn{2}{l}{Define $k$ integer as $k = d + h$} & closure of the integers, renaming $d + h$ \\[6pt]
& ${}= 2(k)$ & substitution \\[6pt]
\end{tabular}

\medskip
$\bullet$ Thus we see that $C = 2(k)$, \quad $k$ is an integer, so by definition of even, $C$ is even.

\medskip
\qed

\sectionbreak
\textbf{Use Warm Up THMs 2--4 for ungraded practice.}
\bigbreak

\textbf{Warm Up THM 2}
\smallbreak
For all positive integers, $A$, $B$, if $A$ is even and $B$ is odd, $A + B = C$ is odd.

\bigbreak

\begin{tabular}{r@{}l@{\hspace{2cm}}l}
\multicolumn{2}{l}{There is $d$ such that $A = 2d$} & definition of even \\[6pt]
\multicolumn{2}{l}{There is $h$ such that $B = 2h + 1$} & definition of odd \\[6pt]
$C$ & ${}= A + B$ & premise \\[6pt]
& ${}= 2d + (2h + 1)$ & substitution \\[6pt]
& ${}= 2d + 2h + 1$ & algebra \\[6pt]
& ${}= 2(d + h) + 1$ & factoring \\[6pt]
\multicolumn{2}{l}{Define $k$ integer as $k = d + h$} & closure of the integers, renaming $d + h$ \\[6pt]
& ${}= 2(k) + 1$ & substitution \\[6pt]
\end{tabular}

\medskip
$\bullet$ Thus we see that $C = 2(k) + 1$, \quad $k$ is an integer, so by definition of odd, $C$ is odd.

\medskip
\qed

\bigbreak
\sectionbreak

\textbf{Warm Up THM 3}
\smallbreak
For all positive integers, $A$, $B$, if $A$ is odd and $B$ is odd, $A + B = C$ is even.
\bigbreak

Let $\mathbb{Z}$ be the set of all integers, then:
\bigbreak
\begin{tabular}{r@{}l@{\hspace{2cm}}l}
\multicolumn{2}{l}{Let $A = 2x + 1$ such that $x \in \mathbb{Z}_{\geq 0}$} & definition of odd \\[6pt]
\multicolumn{2}{l}{Let $B = 2y + 1$ such that $y \in \mathbb{Z}_{\geq 0}$}  & definition of odd \\[6pt]
$C$ & ${}= A + B$ & premise \\[6pt]
& ${}= (2x + 1) + (2y + 1)$ & substitution \\[6pt]
& ${}= 2x + 2y + 2$ & algebra \\[6pt]
& ${}= 2(x + y + 1)$ & factoring \\[6pt]
\multicolumn{2}{l}{Let $z$ be an integer $z = x+y+1$} & closure of integers \\[6pt]
& ${}= 2(z)$ & substitution \\[6pt]

\end{tabular}

\medskip

$\bullet$ Thus we have shown that $C = 2(z)$ such that $z \in \mathbb{Z}_{\geq 0}$, so by definition $C$ is even.

\medskip
\qed

% YOUR PROOF HERE

\sectionbreak
\textbf{Warm Up THM 4}
\smallbreak
Prove for all positive integers, $A$, $B$, if $A$ and $B$ are even, $A \cdot B = C$ is even.

\bigbreak

% YOUR PROOF HERE

\bigbreak
\sectionbreak

\textbf{Warm Up THM 5}
\smallbreak
Prove: For all positive integers, $A$, $B$, if $A$ is even and $B$ is odd, $A \cdot B = C$ is even.

\bigbreak

% YOUR PROOF HERE

\sectionbreak
\textbf{\#1 (Graded) -- Warm Up THM 6}
\medbreak
Prove: For all positive integers, $A$, $B$, if $A$ and $B$ are odd, $A \cdot B = C$ is \textbf{odd}.

\bigbreak
We will now prove that $A \cdot B = C$ is odd.
\bigbreak
\begin{tabular}{r@{}l@{\hspace{2cm}}l}
    \multicolumn{2}{l}{There exists an integer $j$ such that $A = 2j + 1$} & definition of odd \\[6pt]
    \multicolumn{2}{l}{There exists an integer $k$ such that $B = 2k + 1$} & definition of odd \\[6pt]
    $C$ & ${}= A \cdot B$ & premise \\[6pt]
    & ${}= (2j + 1) \cdot (2k + 1)$ & substitution \\[6pt]
    & ${}= 4jk + 2j + 2k + 1$ & FOIL \\[6pt]
    & ${}= 2(2jk + j + k) + 1$ & factoring \\[6pt]
    \multicolumn{2}{l}{Let $h$ be an integer such that $h = 2jk + j + k$} & closure of the integers \\&&(under multiplication and addition)\\[6pt]
    & ${}= 2(h) + 1$ & substitution \\[6pt]
\end{tabular}

\medskip

$\bullet$ Thus we have shown that $C = 2h + 1$, therefore by the definition of an odd integer, $C$ is an odd integer.

\medskip
\qed

\sectionbreak
\textbf{\#2} Prove that if $x$ is odd, then $5x + 3$ is ``\ldots'' two different ways.

\bigbreak

\begin{itemize}
    \item A direct proof by \textbf{applying the Warm-Up Thms} -- this means just USE the RESULTS of what you have proved above.

    \bigbreak
    We will now prove that if $x$ is an odd integer, then $5x + 3$ is \textbf{even}.
    \bigbreak

    \begin{tabular}{r@{}l@{\hspace{2cm}}l}
    \multicolumn{2}{l}{Let $5x = A$, then we know A is an odd integer.} & We deduce this from Theorem 6 \\ && which states if $A$ and $B$ are odd integers, \\ && then $A\cdot B = C$ is also an odd integer.\\[6pt]
    \multicolumn{2}{l}{Let $B$ be an odd integer such that $B = 3$.} & By definition, $B$ is odd \\ && because 3 is an odd number.\\[6pt]
    $5x + 3$ & ${}= A + B$ & by substitution \\[6pt]
    \end{tabular}
    \bigbreak
    $\bullet$ By Theorem 3, the sum of two odd integers is even. Thus, we have shown that if $x$ is odd, then $5x + 3$ is even.
    \medskip
    \qed
    % GRADED -- YOUR PROOF HERE

    \bigbreak

    \item Using the \textit{techniques} of the warm up Thm.\ and the formal definitions of even and odd. (This will look a lot like your WUT proofs.)

    \bigbreak
    We will now prove that if $x$ is an odd integer, then $5x + 3$ is \textbf{even}.
    \bigbreak

    \begin{tabular}{r@{}l@{\hspace{2cm}}l}
    \multicolumn{2}{l}{Let $x = 2m + 1$, where $m \in \mathbb{Z}_{\geq 0}$} & definition of an odd number \\[6pt]
    $5x + 3$ & ${}= 5(2m+1) + 3$ & by substitution \\[6pt]
    & ${}= 10m + 5 + 3$ & distribution/multiplication \\[6pt]
    & ${}= 10m + 8$ & addition \\[6pt]
    & ${}= 2(5m + 4)$ & factoring \\[6pt]
    \multicolumn{2}{l}{\quad Let $n$ be an integer greater than zero such that $n = 5m + 4$} & closure of the integers \\&&(under multiplication and addition)\\[6pt]
    & ${}= 2(n)$ & by substitution
    \end{tabular}
    \bigbreak
    $\bullet$ Thus, $5x + 3 = 2n$ when $n$ is an integer greater than zero. Therefore by definition, $5x + 3$ is even.
    \medskip

    \qed
    % GRADED -- YOUR PROOF HERE

\end{itemize}

\sectionbreak
\textbf{Ungraded reflection:} Explain the difference between using the RESULTS of a Theorem and using the TECHNIQUES of a theorem's proof.

\bigbreak

% YOUR ANSWER HERE

\sectionbreak
\textbf{\Large You may NOT use Warm-Up THM results again unless specified -- this includes exams.}

\bigbreak

\textbf{Ungraded:} Attempt to use a direct proof (similar to techniques of WUT THMs):
\medbreak
\quad If $3n + 7$ is odd, then $n$ is even ($n$ is an integer).

\bigbreak

% YOUR ATTEMPT HERE

\bigbreak

\textit{Why} is this not working? (hint: it should not work)

\bigbreak

% YOUR EXPLANATION HERE

\medbreak
This is a good one to discuss on Piazza.

\sectionbreak
\textbf{\#3} Prove using a proof by contrapositive. (Be sure to use the right format.)
\medbreak
\quad If $3n + 7$ is odd, then $n$ is even ($n$ is an integer).

\bigbreak
We will prove this by contrapositive. Lets assume that $n$ is odd, then we will prove $3n + 7$ is even.
\bigbreak

\begin{tabular}{r@{}l@{\hspace{2cm}}l}
    \multicolumn{2}{l}{Let $n = 2a + 1$ such that $a$ is a positive integer.} & By the definition of an odd number \\[6pt]
    $3n + 7$ & ${}= 3(2a + 1) + 7$ & by substitution \\[6pt]
    & ${}= 6a + 3 + 7$ & by distribution and multiplication \\[6pt]
    & ${}= 6a + 10$ & by addition \\[6pt]
    & ${}= 2(3a + 5)$ & factoring \\[6pt]
    \multicolumn{2}{l}{Let $b = 3a + 5$, then $b$ is an integer} & closure of the integers \\&&(under multiplication and addition) \\[6pt]
    & ${}= 2(b)$ & by substitution\\[6pt]
\end{tabular}

\bigbreak
Thus, we have shown that if $n$ is odd, then $3n + 7$ is even. Then by contrapositive, If $3n + 7$ is odd, then $n$ is even.

\medskip

\qed

\bigbreak

Why does this work, when a direct proof did not?
\medskip

In this case, a direct proof would not allow us to begin with the clean definition of an odd or even integer. This leaves us with a poor setup to proving $n$ is even, and we get stuck in the middle of the proof. Starting with $3n + 7 = 2k + 1$ gives you $n = \frac{2k - 6}{3}$, and dividing by 3 makes it impossible to cleanly factor out a 2 to show $n$ is even.  Since the definition of a direct proof is $p \to q$ where $p$ is the expression "$3n + 7$ is odd" and q is "$n$ is even", a proof by contrapositive ($\neg q \to \neg p$) allows us to flip the expression, which includes the starting point of our proof to "$n$ is odd", allowing us to start with a definition making the proof much easier.

\bigbreak

% GRADED -- YOUR EXPLANATION HERE

\sectionbreak
\textbf{\#4} Prove for positive integers, if $n^2$ is even, then $n$ is even. \quad $\langle$careful, careful, careful$\rangle$

\bigbreak
We will again prove this by contrapositive. Assume $n$ is odd.
\bigbreak
\begin{tabular}{r@{}l@{\hspace{2cm}}l}
    \multicolumn{2}{l}{Let $n = 2a + 1$ such that $a$ is a positive integer.} & By the definition of an odd number \\[6pt]
    $n^2$ & ${}= (2a + 1)^2$ & using substitution \\[6pt]
    & ${}= 4a^2 + 4a + 1$ & FOIL \\[6pt]
    & ${}= 2(2a^2 + 2a) + 1$ & factoring 2 out \\[6pt]
    & ${}= 2(2a(a+1)) + 1$ & factoring $a$ out \\[6pt]
    \multicolumn{2}{l}{Let $b = a+1$ then $b$ is an integer} & closure of the integers \\&&(under addition) \\[6pt]
    \multicolumn{2}{l}{Let $c = 2a(a+1) = 2ab$, then $c$ is an integer} & closure of the integers under multiplication \\ && since we have shown $b$ is an integer \\[6pt]
    & ${}= 2(c) + 1$ & by substitution
\end{tabular}
\bigbreak
Thus we have shown that $n^2 = 2(c) + 1$ and therefore by definition $n^2$ is odd. By contrapositive, since we have shown that if $n$ is odd, then $n^2$ is odd, then if $n^2$ is even, $n$ is also even.
\medskip

\qed

\sectionbreak
\textbf{\#5} Prove that $13n + 3$ is even if and only if $n$ is odd. $n$ is an integer. (Must prove both directions, why?)

\bigbreak
The above statement follows the same logic pattern as a proof of equivalence. Proofs of equivalence follow the tautology $(p \leftrightarrow q) \leftrightarrow (p \to q) \land (q\to p)  $. So, in order to prove the statement above, we must prove both direction $p \to q$ and $q \to p$ where $p$ is "$13n + 3$ is even" and $q$ is "$n$ is odd".
\bigbreak
\textbf{Direction 1:} If $n$ is odd, then $13n + 3$ is even.

\bigbreak
We will prove this statement directly.
\bigbreak

\begin{tabular}{r@{}l@{\hspace{2cm}}l}
    \multicolumn{2}{l}{Let $n = 2q + 1$ such that $q$ is a positive integer.} & By the definition of an odd number \\[6pt]
    $13n + 3$ & ${}= 13(2q + 1) + 3$ & by substitution \\[6pt]
    & ${}= 26q + 13 + 3$ & by distributive multiplication \\[6pt]
    & ${}= 26q + 16$ & by addition \\[6pt]
    & ${}= 2(13q + 8)$ & factoring 2 \\[6pt]
    \multicolumn{2}{l}{Let $p = 13q + 8$ then $p$ is an integer} & closure of the integers \\&&(under addition) \\[6pt]
    & ${}= 2(p)$ & by substitution \\[6pt]
\end{tabular}    
\bigbreak
Thus, we have shown that $13n + 3 = 2p$ and therefore by definition, $13n + 3$ is even.

\bigbreak
\bigbreak

\textbf{Direction 2:} If $13n + 3$ is even, then $n$ is odd.

\bigbreak

We will prove this by contraposition. Assuming $n$ is even, we have
\bigbreak
\begin{tabular}{r@{}l@{\hspace{2cm}}l}
    \multicolumn{2}{l}{Let $n = 2r$ such that $r$ is a positive integer.} & By the definition of an even number \\[6pt]
    $13n + 3$ & ${}= 13(2r) + 3$ & By substitution \\[6pt]
    & ${}= 26r + 3$ & By multiplication \\[6pt]
    & ${}= 26r + (1 + 2)$ & By substitution ($1+2=3$) \\[6pt]
    & ${}= (26r + 2) + 1$ & By associativity \\[6pt]
    & ${}= 2(13r + 1) + 1$ & factoring \\[6pt]
    \multicolumn{2}{l}{Let $s = 13r + 1$ then $s$ is an integer} & closure of the integers \\&&(under addition) \\[6pt]
    & ${}= 2(s) + 1$ & by substitution \\[6pt]
\end{tabular}
\bigbreak
Thus, we have shown that $13n + 3 = 2s + 1$ and is therefore by definition odd. So, since we have shown that if $n$ is even then $13n + 3$ is odd, then by contraposition if $13n + 3$ is even, then $n$ is odd.

\bigbreak

Since we have shown both \textbf{Direction 1} and \textbf{Direction 2} are true, then we have proven the equivalence "$13n + 3$ is even if and only if $n$ is odd".
\medskip

\qed

\sectionbreak
\textbf{Ungraded -- try this in your own words.}
\medbreak
Prove there is no largest integer $z$. Assume the domain of integers.
\medbreak
\textit{There is no integer $z$, for all $n$, $z > n$.}
\medbreak
Use a proof by contradiction. Review the directions for a proof by contradiction. After your proof briefly discuss -- What does this imply about the nature of infinity? Can infinity be an integer?

\bigbreak

\textbf{Say ``We prove by\ldots''}

\bigbreak

For this, we will prove by contradiction.

\bigbreak

\textbf{State the negation --}

\bigbreak

Assume that there exists an integer $z$, such that for all $n$, $z > n$ where $z$ and $n$ are both integers.

\bigbreak

\textbf{Consider the negation.}

\bigbreak
Assuming the above statement to be true, we expect when proving that a step will be logically contradictary.
\bigbreak

\textbf{Do something valid with the negation that shows something goes wrong:}

\bigbreak

If this is true, then:

\begin{tabular}{l@{\hspace{2cm}}l}
    Let $n = z + 1$ & if $z$ is an integer than so is $z + 1$ \\[6pt]
    z > (z + 1) & substitution
\end{tabular}

\bigbreak

\textbf{State ``Contradiction!!!''}

\bigbreak

But we know that by the definition of integers and addition, $z < z + 1$ !

\textbf{State ``by contradiction we have\ldots''}

\bigbreak

Thus, we have a contradiction because by the laws of addition and integers $z > (z + 1)$ must be false.

\bigbreak
\textit{What does this imply about the nature of infinity? Can infinity be an integer?}

\bigbreak

This implies that infinity is not an integer! Which is pretty strange to think about but it makes sense. While infinity does not exist this does make sense in terms of the laws of the set of integers. Essentially if you try to pick a $z$ that is greater than all $n$, you will never find $z$ because there is always an $n$ thats greater, meaning the set is unbounded and infinity does not exist.

\sectionbreak
\textbf{\#6} Assume the domain of positive integers. Prove if $n = ab$, then $a \leq \sqrt{n}$ or $b \leq \sqrt{n}$.
\medbreak
Prove by contradiction (using required steps).

\bigbreak

We will prove this by contradiction. Assume for contradiction that there exists an $n = ab$, $\neg(a \leq \sqrt{n})\equiv a > \sqrt{n}$ and $\neg(b \leq \sqrt{n})\equiv b > \sqrt{n}$ (by DeMorgan's law). Then if $a > \sqrt{n}$ and $b > \sqrt{n}$:
\bigbreak
\begin{tabular}{r@{}l@{\hspace{2cm}}l}
    \multicolumn{2}{l}{Let $a > \sqrt{n}$} & As assumed in our initial statement \\[6pt]
    \multicolumn{2}{l}{Let $b > \sqrt{n}$} & As assumed in our initial statement  \\[6pt]
    $ab$ & ${} > \sqrt{n} \cdot \sqrt{n}$ & In this case because $a, b, n > 0$, by laws of multiplication this statement is true\\[6pt]
    & ${} > n$ & by multiplication laws\\[6pt]
\end{tabular}

\bigskip
 But this leads to a contradiction, because we assumed $n = ab$ but here, $ab > n$! Thus, we have shown there does not exist an $n = ab$ such that $a > \sqrt{n}$ and $b > \sqrt{n}$, and therefore if $n = ab$, then $a \leq \sqrt{n}$ or $b \leq \sqrt{n}$ is true by contradiction.
\bigskip

 \qed

\sectionbreak
\textbf{\#7} We define an exponent $b$ of $a$, ($a$ is an element of the positive real numbers, and $b$ is an element of the natural numbers) as:

\medbreak
$a^b$ is $a$ multiplied by itself $b$ times. Example $a^4 = a \cdot a \cdot a \cdot a$

$a^0 = 1$

\medbreak
We will use only the basic principles of algebra (associativity, commutativity, etc.), and the definition of exponents given above to prove the rules of exponents.

\medbreak
Assume $a$ and $b$ are elements of the positive real numbers. You must explain each step. Notice the format of this example proof and follow this structure for \#8 and \#9. Do not ``work both sides of the equation.''

\medbreak
Notice how this proof does not use the rule it is trying to prove. \textbf{For \#7, just fill in the blanks with the correct justification of each step.}

\bigbreak

Prove that $a^3 \cdot a^5 = a^{(3+5)}$

\bigbreak

\begin{tabular}{r@{}l@{\hspace{2cm}}l}
$a^3 \cdot a^5$ & ${}= (aaa) \cdot (aaaaa)$ & Def.\ of exponents \\[6pt]
& ${}= (aaaaaaaa)$ & \textbf{Associative} law \\[6pt] % GRADED -- FILL IN BLANK
& ${}= a^8$ & Definition of \textbf{exponents} \\[6pt] % GRADED -- FILL IN BLANK
& ${}= a^{(3+5)}$ & Substitution using $8 = 3 + 5$ \\[6pt]
\end{tabular}

\medskip
\qed

\sectionbreak
\textbf{\#8} Prove that $(a^{3})^{5} = a^{3 \cdot 5}$

\bigbreak
We will prove this directly. We have:
\bigbreak
\begin{tabular}{r@{}l@{\hspace{2cm}}l}
    $(a^{3})^{5}$ & ${}= (aaa)^{5}$ & definition of exponents \\[6pt]
    & ${}= (aaa)\cdot(aaa)\cdot(aaa)\cdot(aaa)\cdot(aaa)$ & definition of exponents (again) \\[6pt]
    & ${}= aaaaaaaaaaaaaaa$ & associative law of multiplication \\[6pt]
    & ${}= a^{15}$ & Definition of exponents \\[6pt]
    & ${}= a^{3 \cdot 5}$ & Substitution using $15 = 3 \cdot 5$ \\[6pt]
\end{tabular}
\bigbreak

\qed

\sectionbreak
\textbf{\#9} Prove that $a^5 \cdot b^5 = (ab)^5$

\bigbreak
We will prove this statement directly. So, we have:
\bigbreak
\begin{tabular}{r@{}l@{\hspace{2cm}}l}
    $(ab)^5$ & ${}= (aaaaa)\cdot b^5$ & definition of an exponent (on a)\\[6pt]
    & ${}= (aaaaa)\cdot (bbbbb)$ & definition of an exponent (on b)\\[6pt]
    & ${}= aaaaabbbbb$ & by multiplication \\[6pt]
    & ${}= ababababab$ & by the law of commutativity \\[6pt]
    & ${}= (ab) \cdot (ab) \cdot(ab) \cdot (ab) \cdot (ab) $ & by the law of associativity \\[6pt]
    & ${}= (ab)^5$  & definition of an exponent \\[6pt]
\end{tabular}
\bigbreak

\qed

\sectionbreak
\textbf{\#10} Considering the above, when we add the definition of exponents to the algebra, are the ``rules'' of exponents really a new idea or just a consequence of the existing structure of algebra? Give a thoughtful and complete answer.

\medskip

Rules of exponents for what we have covered so far are just a consequence of the existing structure of algebra, or even just shorthand when you need to show a mathematical expression multiplied a given number of times. Generally, from the proofs on \#7, \#8, and \#9, all of the steps boil down to multiplication laws, associativity, and commutativity. I do feel this only holds true if you have $a^n$ where $n$ is a natural number or positive integer. When looking at $a^n$ when $n$ becomes negative, rational, or even irrational, the idea of an exponent and proofs with exponents get much more complex and nontrivial. However, with the help of expanded theorems and the definition of logarithms, proofs of more complex exponent statements become more feasible versus only working with basic mathematical and logical axioms.

\bigbreak

% GRADED -- YOUR ANSWER HERE

\sectionbreak
\textbf{Notes} -- ungraded, add as many pages as you like.
\medbreak
Use these questions or other notes to help you organize your study and prepare for the exam.

\begin{enumerate}
    \item Which ideas were new to you?
    \item What questions might make good exam questions?
    \item What ideas will you want to review later?
    \item What were some ``a-ah'' moments, or big ideas you want to remember?
\end{enumerate}

\bigbreak

% YOUR NOTES HERE

\end{document}
